\section{Open Questions}

The following five questions arise naturally from this replication and remain genuinely open. Each includes a concrete next step that could be executed with the current PennyLane replication as a starting point.

\subsection*{Q1. Scaling to larger molecules and the measurement-budget wall}
\textbf{Question.} How does the Peruzzo-style VQE scale when extended from HeH$^+$ (2 qubits, STO-3G) to mid-size molecules like BeH$_2$, H$_2$O, or LiH in a larger basis (6-31G, cc-pVDZ), where the Pauli term count grows as $O(N^4)$ and Hamiltonian averaging becomes measurement-budget-limited?

\textbf{Basis.} The paper demonstrates on a single 2-qubit molecule; the algorithmic core is measurement-count-bounded, and modern VQE literature (McClean et al.\ 2016, Kandala et al.\ 2017) shows the shot budget dominates practical runtime for larger systems.

\textbf{Next steps.} Extend \verb|replicate.py| to BeH$_2$/H$_2$O/LiH under STO-3G and 6-31G. Report Pauli term counts, per-iteration measurement budgets under a shot-noise model, and wall-clock convergence vs.\ FCI. Compare naive term-by-term averaging to grouped-Pauli measurements (qubit-wise commuting, general commuting). Target: quantify the measurement-count blowup vs.\ molecule size and basis.

\subsection*{Q2. Modern ansatze: ADAPT-VQE and k-UpCCGSD}
\textbf{Question.} Does ADAPT-VQE or k-UpCCGSD (adaptive/compact ansatze developed after 2014) recover the HeH$^+$ dissociation curve with strictly fewer parameters and a shallower circuit than the fixed UCCSD used here, without loss of chemical accuracy?

\textbf{Basis.} UCCSD is fixed-depth and includes all single/double excitations regardless of whether they contribute; ADAPT-VQE (Grimsley et al.\ 2019) grows the ansatz from an operator pool by largest energy gradient, and k-UpCCGSD (Lee et al.\ 2019) is a compact generalized-singles-doubles variant. Both post-date the paper.

\textbf{Next steps.} Add ADAPT-VQE and k-UpCCGSD ansatze alongside UCCSD in \verb|replicate.py|. On the HeH$^+$ dissociation curve, report (parameter count, CNOT depth, VQE$-$FCI error) triples per bond length for each ansatz. Compare against the fixed UCCSD baseline. Then repeat on BeH$_2$ where the compression is expected to be larger.

\subsection*{Q3. Photochemistry benchmarks and static correlation}
\textbf{Question.} How does this replication perform on a photochemistry benchmark set (e.g., the singlet--triplet gaps of small carbenes, or the QMSpin / QM7-X excitation subset) that stresses static correlation and near-degeneracies, which the 2-qubit HeH$^+$ demonstration cannot expose?

\textbf{Basis.} HeH$^+$ is a closed-shell 2-electron system where UCCSD is exact in the minimal basis; the paper's demonstration therefore does not test the regimes where VQE variants diverge in performance (multireference character, spin contamination, strong static correlation).

\textbf{Next steps.} Select a small photochemistry benchmark (methylene singlet--triplet gap; ozone equilibrium bond lengths; a QMSpin subset). For each, run VQE with UCCSD and ADAPT-VQE against FCI/CASSCF baselines. Report (chem.\ accuracy achieved?, ansatz depth, whether the correct spin state was reached). Explicitly note failure modes if the ansatz cannot reach the reference.

\subsection*{Q4. Realistic gate-noise robustness}
\textbf{Question.} Under a realistic gate-noise model (depolarizing + amplitude damping at superconducting-transmon error rates, $\sim 10^{-3}$ per two-qubit gate), does the noiseless HeH$^+$ chemical-accuracy result survive, or does it collapse toward the paper's 96\% hardware figure?

\textbf{Basis.} The paper's 96\%-within-chemical-accuracy number is a hardware statement, and no noise was included in this replication. Modern noise-injection tooling (PennyLane \verb|default.mixed|, Qiskit noise models) makes a controlled sweep straightforward.

\textbf{Next steps.} Add a noise sweep to \verb|replicate.py| using PennyLane \verb|default.mixed| with two-qubit depolarizing rates from $10^{-4}$ to $10^{-2}$. For each rate, compute (mean VQE$-$FCI error across the HeH$^+$ curve, chem-accuracy fraction). Plot chem-accuracy fraction vs.\ gate error rate and identify the rate at which the noiseless 100\% falls to the paper-reported 96\%. Cross-check by using published IBM Falcon/Eagle two-qubit error rates as the anchor.

\subsection*{Q5. Photonic vs.\ superconducting-transmon: post-2014 platform comparison}
\textbf{Question.} How do the runtime and error trends of superconducting-transmon VQE demonstrations (2018--2024, e.g., Kandala et al., Google Sycamore, IBM heavy-hex) compare to the 2014 photonic demonstration on the same HeH$^+$ target, and does the photonic platform still have any advantage (photon-loss vs.\ decoherence, native connectivity) for VQE-class algorithms?

\textbf{Basis.} The paper is the photonic-processor VQE demonstration; the field has since shifted almost entirely to superconducting-transmon hardware. A like-for-like HeH$^+$ comparison across platforms has not, to the replicator's knowledge, been systematically compiled.

\textbf{Next steps.} Literature review: enumerate all published HeH$^+$ or H$_2$ VQE hardware runs 2014--2024. Extract (platform, qubit count, two-qubit error rate, best VQE$-$FCI mHa, runtime, ansatz). Tabulate. Then answer: on the same molecule, has any post-2014 hardware beaten the 2014 photonic chip on (a) accuracy, (b) time-to-solution, (c) shot count? Identify the platform-agnostic vs.\ platform-specific parts of any advantage.
